% ============================================================================
% DÉFINITION DU COURS 04 : CINÉMATIQUE
% ============================================================================

\titlehead{Thomas Lusseau - SII en ATS}
\title[Cinématique]{Détermination des lois de mouvement — Cinématique}
\author[TL]{Thomas Lusseau}
\date{\today}

\TLCourseCoverSetup
  {\repStyle/FondPoly.png}
  {images/image5.jpeg}
  {Cycle 2 : Modéliser et déterminer les lois de\\mouvement d'un mécanisme}
  {Détermination des lois de mouvement\\Approche cinématique}
  {Lycée Robert Doisneau - ATS}
  {Thomas Lusseau}

\TLCourseCoverContents{%
  \TLCourseContentsLine{4.1}{Objectifs et introduction}{sec:cine-objectifs}
  \TLCourseContentsLine{4.2}{Produit vectoriel et torseurs}{sec:cine-produit}
  \TLCourseContentsLine{4.3}{Mouvements et trajectoires}{sec:cine-mouvements}
  \TLCourseContentsLine{4.4}{Position, vitesse et accélération}{sec:cine-pva}
  \TLCourseContentsLine{4.5}{Dérivation et composition}{sec:cine-derivation}
  \TLCourseContentsLine{4.6}{Contact et roulement sans glissement}{sec:cine-contact}
  \TLCourseContentsLine{4.7}{Loi de vitesse en trapèze}{sec:cine-trapeze}
  \TLCourseContentsLine{4.8}{Exercices du chapitre}{sec:cine-exercices}
}

\TLCourseFooter
\TLCourseCover
\markboth{Cinématique}{Cinématique}

\livrettrue
\collefalse

\definecolor{TLCineBlue}{RGB}{0,84,127}
\definecolor{TLCineGreen}{RGB}{72,145,77}
\definecolor{TLCineOrange}{RGB}{214,125,25}

\newtcolorbox{TLCineMethod}[1]{
  enhanced,breakable,colback=TLCineBlue!6,colframe=TLCineBlue,
  colbacktitle=TLCineBlue,coltitle=white,title={Méthode — #1},
  fonttitle=\bfseries,boxrule=.8pt,arc=1.5mm,
  left=4mm,right=4mm,top=3mm,bottom=3mm,
  before skip=.9\baselineskip,after skip=.9\baselineskip
}
\newtcolorbox{TLCineDef}[1]{
  enhanced,breakable,colback=TLCineGreen!6,colframe=TLCineGreen!80!black,
  colbacktitle=TLCineGreen!80!black,coltitle=white,title={Définition — #1},
  fonttitle=\bfseries,boxrule=.8pt,arc=1.5mm,
  left=4mm,right=4mm,top=3mm,bottom=3mm,
  before skip=.9\baselineskip,after skip=.9\baselineskip
}
\newtcolorbox{TLCineApp}[1]{
  enhanced,breakable,colback=TLCineOrange!7,colframe=TLCineOrange,
  colbacktitle=TLCineOrange,coltitle=white,title={Application — #1},
  fonttitle=\bfseries,boxrule=.8pt,arc=1.5mm,
  left=4mm,right=4mm,top=3mm,bottom=3mm,
  before skip=.9\baselineskip,after skip=.9\baselineskip
}

\newcommand{\TLCineImage}[2][.78\linewidth]{%
  \IfFileExists{#2}{\includegraphics[width=#1,keepaspectratio]{#2}}{}%
}

\TLStartCourseChapter{4}{Détermination des lois de mouvement — Cinématique}{chap:cinematique}

% ============================================================================
% COURS 04 — CINÉMATIQUE — PARTIE A
% Source : 5-Cinematique2022.docx
% ============================================================================

\section{Objectifs et mises en situation}
\label{sec:cine-objectifs}

\begin{TLCineMethod}{Objectifs du cours}
Décrire et caractériser les mouvements de tous les points des solides
d'un système : trajectoire, position, vitesse et accélération, sans
étudier les causes mécaniques qui produisent ces mouvements.
\end{TLCineMethod}

\begin{center}
\small
\begin{tabularx}{\linewidth}{>{\bfseries}p{2cm}X c}
\toprule
Je connais & Connaissances attendues & \\
\midrule
& Les définitions des vecteurs position, vitesse et accélération. & $\square$\\
& Les relations de composition des vitesses, rotations et torseurs cinématiques. & $\square$\\
& Le champ des vecteurs vitesse d'un solide indéformable. & $\square$\\
& La loi de vitesse en trapèze. & $\square$\\
\midrule
Je sais & Nommer un mouvement et décrire une trajectoire. & $\square$\\
& Dériver un vecteur dans un repère mobile. & $\square$\\
& Déterminer une vitesse par dérivation ou composition. & $\square$\\
& Déterminer une accélération par dérivation du vecteur vitesse. & $\square$\\
& Passer d'un profil d'accélération à la vitesse et à la position. & $\square$\\
\bottomrule
\end{tabularx}
\end{center}

\subsection{Robot de chirurgie}

La chirurgie robotique impose de maîtriser les vitesses et les
trajectoires. Une mauvaise préparation ou un maniement hésitant peut
allonger l'intervention. Le concepteur doit donc relier les ordres
envoyés aux actionneurs à la cinématique réellement imposée aux
effecteurs.

\begin{center}
\TLCineImage[.42\linewidth]{images/image5.jpeg}\hfill
\TLCineImage[.48\linewidth]{images/image6.png}
\end{center}

\subsection{Prothèse transtibiale}

Une prothèse doit reproduire les amplitudes, vitesses et accélérations
d'un pied réel. La modélisation cinématique permet de vérifier les
mouvements possibles et les niveaux d'accélération transmis à
l'utilisateur.

\TLCineImage[.48\linewidth]{images/image7.jpeg}

\section{Introduction à la cinématique}
\label{sec:cine-introduction}

La cinématique est la partie de la mécanique qui décrit le mouvement,
indépendamment des efforts qui le produisent. Les solides sont supposés
\textbf{indéformables} : pour deux points $A$ et $B$ d'un même solide,
\[
\forall t,\qquad \|\overrightarrow{AB}\|=\mathrm{constante}.
\]

\subsection{Base, repère et référentiel}

Une base orthonormée directe est notée
$\mathcal B=(\vec x,\vec y,\vec z)$. Un repère d'espace associe cette
base à une origine $O$ :
\[
\mathcal R=(O,\vec x,\vec y,\vec z).
\]
Un référentiel comprend un repère d'espace et un repère de temps. En
SII, le temps est commun à tous les référentiels ; on parle donc souvent
simplement de repère de référence.

Dans la suite, un solide et le repère qui lui est lié sont souvent
confondus dans les notations.

\section{Produit vectoriel et produit mixte}
\label{sec:cine-produit}

\subsection{Produit vectoriel}

Le produit vectoriel $\vec U\wedge\vec V$ est perpendiculaire au plan
$(\vec U,\vec V)$, orienté de sorte que
$(\vec U,\vec V,\vec U\wedge\vec V)$ soit direct, et de norme :
\[
\boxed{\|\vec U\wedge\vec V\|
=\|\vec U\|\,\|\vec V\|\,|\sin(\vec U,\vec V)|.}
\]

Dans une base orthonormée directe :
\[
\vec x\wedge\vec y=\vec z,
\quad
\vec y\wedge\vec z=\vec x,
\quad
\vec z\wedge\vec x=\vec y.
\]
L'inversion de l'ordre change le signe.

\begin{TLCineDef}{Propriétés}
\[
\vec U\wedge\vec V=-\vec V\wedge\vec U,
\]
\[
\vec U\wedge(\vec V+\vec W)
=\vec U\wedge\vec V+\vec U\wedge\vec W,
\]
\[
(\lambda\vec U)\wedge(\mu\vec V)
=\lambda\mu(\vec U\wedge\vec V).
\]
Le produit vectoriel est nul lorsque les deux vecteurs sont colinéaires
ou lorsque l'un d'eux est nul.
\end{TLCineDef}

Dans une base orthonormée :
\[
\vec U\wedge\vec V=
\begin{vmatrix}
\vec x&\vec y&\vec z\\
U_x&U_y&U_z\\
V_x&V_y&V_z
\end{vmatrix}.
\]

\begin{TLCineApp}{Calcul de produit vectoriel}
Avec
$\vec V_1=2\vec x+2\vec y-3\vec z$ et
$\vec V_2=-\vec x+6\vec z$, calculer
$\vec V_1\wedge\vec V_2$ puis vérifier son orthogonalité avec les deux
vecteurs de départ.
\end{TLCineApp}

\subsection{Produit mixte}

Le produit mixte de trois vecteurs est le réel :
\[
(\vec U,\vec V,\vec W)
=(\vec U\wedge\vec V)\cdot\vec W.
\]
Sa valeur absolue représente le volume du parallélépipède construit sur
les trois vecteurs. Il est invariant par permutation circulaire et
change de signe lorsqu'on échange deux opérandes.

\section{Torseur cinématique}
\label{sec:cine-torseur}

Le mouvement du solide 1 par rapport au solide 0 est décrit, au point
$A$, par le torseur cinématique :
\[
\boxed{
\left\{\mathcal V_{1/0}\right\}_A
=\left\{\begin{array}{c}
\vec\Omega_{1/0}\\
\vec V(A\in1/0)
\end{array}\right\}_A.}
\]
La résultante $\vec\Omega_{1/0}$ est indépendante du point d'expression.
Le moment cinématique est la vitesse du point considéré.

\begin{TLCineDef}{Inversion du mouvement}
\[
\boxed{\left\{\mathcal V_{2/1}\right\}
=-\left\{\mathcal V_{1/2}\right\}.}
\]
Le point d'expression et la base doivent être précisés pour exploiter
les composantes.
\end{TLCineDef}

\subsection{Changer le point d'expression}

Pour un torseur de résultante $\vec R$ et de moment $\vec M_A$ :
\[
\boxed{\vec M_B=\vec M_A+\overrightarrow{BA}\wedge\vec R.}
\]
Pour le torseur cinématique :
\[
\boxed{\vec V(B\in1/0)=\vec V(A\in1/0)
+\overrightarrow{BA}\wedge\vec\Omega_{1/0}.}
\]
Cette relation constitue le \textbf{champ des vecteurs vitesse} d'un
solide indéformable.

\begin{TLCineApp}{Hélice d'avion}
Une hélice tourne autour de l'axe $(K,\vec w)$ avec
$\vec\Omega_{2/1}=-31\vec w\ \mathrm{rad\,s^{-1}}$. Le point $P$ est
tel que $\overrightarrow{PK}=l\vec u$.
Déterminer $\vec V(P\in2/1)$ en déplaçant le torseur de $K$ vers $P$.
\end{TLCineApp}

\subsection{Opérations sur les torseurs}

Deux torseurs sont égaux si leurs résultantes sont égales et si leurs
moments, exprimés au même point, sont égaux. Pour additionner des
torseurs, il faut les réduire au même point et exprimer les composantes
dans une base commune.

Le comoment de deux torseurs est le réel :
\[
\left\{\mathcal T_1\right\}\otimes
\left\{\mathcal T_2\right\}
=\vec R_1\cdot\vec M_{2,A}
+\vec R_2\cdot\vec M_{1,A}.
\]
Il est indépendant du point de réduction.

% ============================================================================
% COURS 04 — CINÉMATIQUE — PARTIE B
% Source : 5-Cinematique2022.docx
% ============================================================================

\section{Mouvements et trajectoires}
\label{sec:cine-mouvements}

Un mouvement est le déplacement relatif d'un solide observé par rapport
à un solide de référence au cours du temps. La notation $S/0$ désigne
le mouvement du solide $S$ par rapport au solide 0.

\begin{TLCineDef}{Trajectoire}
La trajectoire $\mathcal T_{M\in S/0}$ est le lieu des positions
successives occupées par le point $M$ appartenant au solide $S$, dans
le repère lié au solide 0.
\end{TLCineDef}

Un point de l'espace peut être considéré comme appartenant à un solide
même si aucune matière n'est présente à cet endroit, par exemple un
point situé sur l'axe d'un cylindre creux.

\subsection{Mouvement de translation}

Un solide 2 est en translation par rapport à 1 si :
\[
\boxed{\vec\Omega_{2/1}=\vec0.}
\]
Tous ses points possèdent alors le même vecteur vitesse :
\[
\forall A,B\in2,\qquad
\vec V(A\in2/1)=\vec V(B\in2/1).
\]
On distingue :
\begin{itemize}
  \item la translation rectiligne : trajectoires rectilignes parallèles ;
  \item la translation circulaire : trajectoires circulaires de même rayon ;
  \item la translation curviligne : trajectoires quelconques identiques par translation.
\end{itemize}

\subsection{Mouvement de rotation autour d'un axe fixe}

Le solide 2 est en rotation autour de l'axe $\Delta=(C,\vec u)$ par
rapport à 1 si :
\[
\left\{\mathcal V_{2/1}\right\}_C
=\left\{\begin{array}{c}
\omega\vec u\\\vec0
\end{array}\right\}_C.
\]
Tous les points de l'axe ont une vitesse nulle. Tout point $M$ extérieur
à l'axe décrit un cercle et :
\[
\boxed{\vec V(M\in2/1)=\overrightarrow{MC}\wedge\vec\Omega_{2/1}.}
\]
La norme vaut $V=R|\omega|$, où $R$ est la distance du point à l'axe.

\subsection{Mouvement quelconque}

Lorsqu'un mouvement n'est ni une translation ni une rotation autour d'un
axe fixe, il est qualifié de quelconque. Il peut être analysé comme une
composition de mouvements simples associés aux liaisons successives.

\section{Vecteur rotation et composition}
\label{sec:cine-rotation}

Le vecteur rotation $\vec\Omega_{i/j}$ indique :
\begin{itemize}
  \item la direction de l'axe instantané de rotation ;
  \item la vitesse angulaire en $\mathrm{rad\,s^{-1}}$ ;
  \item le sens positif selon la règle du tire-bouchon.
\end{itemize}

Pour une rotation paramétrée par $\alpha$ autour de $\vec x_1=\vec x_2$ :
\[
\vec\Omega_{2/1}=\dot\alpha\vec x_1.
\]

La composition des rotations s'écrit :
\[
\boxed{\vec\Omega_{n/0}
=\vec\Omega_{n/n-1}+\cdots+\vec\Omega_{2/1}+\vec\Omega_{1/0}.}
\]

\begin{TLCineApp}{Machine cinq axes}
Les solides 1, 4 et 5 sont en translation. Le solide 2 tourne autour de
$\vec x_2$ avec l'angle $\alpha$, et le solide 3 autour de $\vec z_2$
avec l'angle $\beta$.
\[
\vec\Omega_{3/0}
=\dot\beta\vec z_2+\dot\alpha\vec x_2.
\]
Identifier les bases égales et tracer les deux figures de changement de
base.
\end{TLCineApp}

\section{Position, vitesse et accélération}
\label{sec:cine-pva}

\subsection{Vecteur position}

Le vecteur position de $M$ dans le repère
$\mathcal R_0=(O_0,\vec x_0,\vec y_0,\vec z_0)$ est :
\[
\overrightarrow{O_0M}
=x(t)\vec x_0+y(t)\vec y_0+z(t)\vec z_0.
\]
Selon la géométrie, les coordonnées cylindriques ou sphériques peuvent
être plus adaptées.

En coordonnées cylindriques :
\[
\overrightarrow{OM}=r\vec u+z\vec z,
\qquad
\vec u=\cos\theta\vec x+\sin\theta\vec y.
\]
En coordonnées sphériques :
\[
\overrightarrow{OM}=r\vec n,
\qquad
\vec n=\cos\varphi\vec z+\sin\varphi\vec u.
\]

\subsection{Vecteur vitesse}

Le vecteur vitesse est la dérivée du vecteur position dans le repère de
référence :
\[
\boxed{\vec V(M\in S/0)
=\left[\frac{\mathrm d\overrightarrow{O_0M}}{\mathrm dt}\right]_0.}
\]
Il est tangent à la trajectoire et s'exprime en
$\mathrm{m\,s^{-1}}$.

Plus généralement, pour un point fixe $Q$ dans le repère 0 :
\[
\vec V(M\in S/0)
=\left[\frac{\mathrm d\overrightarrow{QM}}{\mathrm dt}\right]_0.
\]

\subsection{Vecteur accélération}

Le vecteur accélération est la dérivée du vecteur vitesse dans le repère
de référence :
\[
\boxed{\vec A(M\in S/0)
=\left[\frac{\mathrm d\vec V(M\in S/0)}{\mathrm dt}\right]_0.}
\]
Il s'exprime en $\mathrm{m\,s^{-2}}$.

\begin{TLCineApp}{Axes de translation d'une machine-outil}
Pour l'extrémité d'outil $P$ :
\[
\overrightarrow{O_0P}
=x(t)\vec x_0+e\vec y_0+(d-z(t)-f)\vec z_0.
\]
On obtient directement :
\[
\vec V(P\in5/0)=\dot x\vec x_0-\dot z\vec z_0,
\]
\[
\vec A(P\in5/0)=\ddot x\vec x_0-\ddot z\vec z_0.
\]
La vitesse maximale vaut
$\sqrt{\dot x_{\max}^2+\dot z_{\max}^2}$.
\end{TLCineApp}

\section{Dérivation vectorielle}
\label{sec:cine-derivation}

Lorsqu'un vecteur est exprimé dans une base mobile, il ne faut pas le
projeter systématiquement dans la base fixe avant de dériver. On utilise
la formule de dérivation vectorielle, dite formule de Boor :

\begin{TLCineDef}{Formule de Boor}
Pour un vecteur $\vec U$ exprimé dans le repère $i$ et dérivé dans le
repère $j$ :
\[
\boxed{
\left[\frac{\mathrm d\vec U}{\mathrm dt}\right]_j
=\left[\frac{\mathrm d\vec U}{\mathrm dt}\right]_i
+\vec\Omega_{i/j}\wedge\vec U.}
\]
Si $\vec U$ est un vecteur unitaire fixe dans $i$ :
\[
\left[\frac{\mathrm d\vec U}{\mathrm dt}\right]_j
=\vec\Omega_{i/j}\wedge\vec U.
\]
\end{TLCineDef}

Pour un repère tournant autour de $\vec x$ avec l'angle $\alpha$ :
\[
\left[\frac{\mathrm d\vec y_2}{\mathrm dt}\right]_0
=\dot\alpha\vec x_2\wedge\vec y_2
=\dot\alpha\vec z_2,
\]
\[
\left[\frac{\mathrm d\vec z_2}{\mathrm dt}\right]_0
=\dot\alpha\vec x_2\wedge\vec z_2
=-\dot\alpha\vec y_2.
\]

\begin{TLCineApp}{Point d'une pièce orientable}
Si
$\overrightarrow{O_0O_3}=y(t)\vec y_0+a\vec z_0+b\vec x_0+c\vec z_2$,
alors :
\[
\vec V(O_3\in3/0)
=\dot y\vec y_0-c\dot\alpha\vec y_2.
\]
En dérivant à nouveau :
\[
\vec A(O_3\in3/0)
=\ddot y\vec y_0-c\ddot\alpha\vec y_2
-c\dot\alpha^2\vec z_2.
\]
Le dernier terme est une accélération centripète.
\end{TLCineApp}

% ============================================================================
% COURS 04 — CINÉMATIQUE — PARTIE C
% Source : 5-Cinematique2022.docx
% ============================================================================

\section{Composition des mouvements}
\label{sec:cine-composition}

\subsection{Composition des vecteurs vitesse}

Pour un point $P$ appartenant au solide 2, lui-même en mouvement par
rapport au solide 1, en mouvement par rapport à 0 :
\[
\boxed{
\vec V(P\in2/0)
=\vec V(P\in2/1)+\vec V(P\in1/0).}
\]
Le premier terme est la \textbf{vitesse relative}; le second est la
\textbf{vitesse d'entraînement}, c'est-à-dire la vitesse du point $P$
considéré comme fixé au solide 1.

Pour $n$ solides successifs :
\[
\vec V(P\in n/0)
=\vec V(P\in n/n-1)+\cdots+\vec V(P\in1/0).
\]
Cette écriture est valable au même point géométrique $P$.

\subsection{Composition des torseurs cinématiques}

\[
\boxed{
\left\{\mathcal V_{n/0}\right\}_P
=\left\{\mathcal V_{n/n-1}\right\}_P+
\cdots+
\left\{\mathcal V_{1/0}\right\}_P.}
\]
Tous les torseurs doivent impérativement être exprimés au même point et
dans une base commune avant addition.

\begin{TLCineMethod}{Déterminer une vitesse par composition}
\begin{enumerate}
  \item Choisir une chaîne de solides reliant le solide étudié au bâti.
  \item Écrire la composition au point demandé.
  \item Pour chaque mouvement simple, déplacer le torseur vers ce point.
  \item Additionner les vecteurs rotation et les vecteurs vitesse.
  \item Laisser si possible le résultat dans les bases naturelles des mouvements.
  \item Vérifier les unités et la direction sur une position particulière.
\end{enumerate}
\end{TLCineMethod}

\subsection{Champ des vecteurs vitesse}

Pour deux points $A$ et $B$ d'un même solide 1 :
\[
\boxed{
\vec V(B\in1/0)
=\vec V(A\in1/0)+\overrightarrow{BA}\wedge\vec\Omega_{1/0}.}
\]
Cette relation ne doit jamais être utilisée entre deux points appartenant
à des solides différents.

\begin{TLCineApp}{Rotation autour d'un axe}
Si $A$ appartient à l'axe de rotation, alors
$\vec V(A\in1/0)=\vec0$ et :
\[
\vec V(B\in1/0)=\overrightarrow{BA}\wedge\vec\Omega_{1/0}.
\]
Le vecteur est tangent au cercle décrit par $B$, et sa norme est
$AB\,\Omega$ lorsque $AB$ est perpendiculaire à l'axe.
\end{TLCineApp}

\subsection{Accélération d'un point d'un solide}

La méthode recommandée consiste à dériver le vecteur vitesse dans le
repère de référence. Pour deux points d'un même solide :
\[
\vec A(B\in1/0)=\vec A(A\in1/0)
+\overrightarrow{BA}\wedge\dot{\vec\Omega}_{1/0}
+(\vec\Omega_{1/0}\wedge\overrightarrow{BA})
\wedge\vec\Omega_{1/0}.
\]
Le deuxième terme est tangentiel et le troisième centripète.

Pour un point appartenant à un solide mobile 2 dans un repère mobile 1 :
\[
\boxed{
\vec A(P\in2/0)=\vec A(P\in2/1)+\vec A(P\in1/0)
+2\vec\Omega_{1/0}\wedge\vec V(P\in2/1).}
\]
Le dernier terme est l'accélération de Coriolis.

\section{Fermeture cinématique}
\label{sec:cine-fermeture}

Dans une boucle cinématique fermée :
\[
\boxed{
\left\{\mathcal V_{1/0}\right\}
+\left\{\mathcal V_{2/1}\right\}
+\cdots+
\left\{\mathcal V_{0/n}\right\}
=\left\{0\right\}.}
\]
Après réduction au même point, les six composantes fournissent un
système d'équations. Cette fermeture permet d'obtenir une loi
entrée-sortie en vitesse, d'étudier les mobilités et d'évaluer le degré
d'hyperstatisme.

\begin{TLCineApp}{Vérin électrique}
Une vis et un écrou sont liés par une liaison hélicoïdale et guidés par
une pivot et une glissière par rapport au bâti. Écrire la fermeture :
\[
\left\{\mathcal V_{2/0}\right\}
-\left\{\mathcal V_{1/0}\right\}
-\left\{\mathcal V_{2/1}\right\}=\{0\}.
\]
Projeter les composantes utiles pour retrouver
$v=(p/2\pi)\omega$ et identifier les équations supplémentaires
traduisant les contraintes internes.
\end{TLCineApp}

\section{Cinématique du contact ponctuel}
\label{sec:cine-contact}

Lorsque deux solides 1 et 2 sont en contact au point géométrique $I$, il
faut distinguer :
\begin{itemize}
  \item le point matériel $I_1$ appartenant au solide 1 ;
  \item le point matériel $I_2$ appartenant au solide 2 ;
  \item le point géométrique de contact $I$.
\end{itemize}

La vitesse de glissement de 2 par rapport à 1 est :
\[
\boxed{
\vec V(I\in2/1)
=\vec V(I\in2/0)-\vec V(I\in1/0).}
\]
Elle appartient au plan tangent commun aux deux solides.

\begin{TLCineDef}{Roulement sans glissement}
Il y a roulement sans glissement au point $I$ lorsque :
\[
\boxed{\vec V(I\in2/1)=\vec0.}
\]
Le point $I$ est alors un centre instantané de rotation pour un
mouvement plan. Cette condition est utilisée pour les engrenages,
roulements, roues sur le sol, poulies-courroies et pignons-chaînes.
\end{TLCineDef}

\begin{TLCineApp}{Roue sur un plateau}
Une roue 2 entraînée par rapport au bras 3 roule sur un plateau 1,
lui-même en rotation par rapport au bâti 0. Écrire :
\[
\vec V(I\in2/1)=\vec V(I\in2/0)-\vec V(I\in1/0)=\vec0.
\]
Décomposer ensuite chacune des vitesses par composition pour relier les
trois vitesses de rotation.
\end{TLCineApp}

\section{Loi de vitesse en trapèze}
\label{sec:cine-trapeze}

Un actionneur ne peut pas passer instantanément d'une vitesse nulle à
une vitesse non nulle. Une loi simple de commande comporte :
\begin{enumerate}
  \item une phase d'accélération constante ;
  \item une phase à vitesse constante ;
  \item une phase de décélération constante.
\end{enumerate}

\begin{center}
\begin{tikzpicture}[>=Latex,x=1.35cm,y=1.2cm]
  \draw[->] (0,0)--(7.2,0) node[right]{$t$};
  \draw[->] (0,0)--(0,3.2) node[above]{$v(t)$};
  \draw[very thick,TLCineBlue]
    (0,0)--(1.6,2.5)--(4.8,2.5)--(6.5,0);
  \draw[dashed] (1.6,0)--(1.6,2.5);
  \draw[dashed] (4.8,0)--(4.8,2.5);
  \draw[dashed] (6.5,0)--(6.5,2.5);
  \node[below] at (1.6,0) {$t_1$};
  \node[below] at (4.8,0) {$t_2$};
  \node[below] at (6.5,0) {$t_3$};
  \node[left] at (0,2.5) {$V$};
\end{tikzpicture}
\end{center}

Les relations fondamentales sont :
\[
v(t)=\dot x(t),
\qquad
a(t)=\dot v(t)=\ddot x(t),
\]
\[
\boxed{x(t)-x(0)=\int_0^t v(T)\,\mathrm dT.}
\]
Le déplacement est donc l'aire algébrique sous la courbe de vitesse.
L'accélération est la pente de cette courbe.

Pour une vitesse maximale $V$, avec accélération de durée $t_a$, palier
de durée $t_p$ et décélération de durée $t_d$, le déplacement total est :
\[
\boxed{
\Delta x=\frac12Vt_a+Vt_p+\frac12Vt_d.}
\]
Si $t_a=t_d$ :
\[
\Delta x=V(t_p+t_a).
\]

\begin{TLCineMethod}{Exploiter un profil de vitesse}
\begin{enumerate}
  \item Lire les instants de changement de phase et les vitesses.
  \item Calculer les accélérations avec $a=\Delta v/\Delta t$.
  \item Déterminer les déplacements par les aires : triangles, rectangles et trapèzes.
  \item Construire la courbe de position, continue et de pente égale à la vitesse.
  \item Vérifier la continuité de la vitesse et de la position aux changements de phase.
\end{enumerate}
\end{TLCineMethod}

\begin{TLCineApp}{Hayon électrique}
Le hayon atteint $20^\circ\,\mathrm{s^{-1}}$ en $0{,}5\ \mathrm s$,
conserve cette vitesse puis décélère symétriquement. Déterminer la durée
du palier pour une ouverture totale de $90^\circ$, puis tracer
$\gamma(t)$ et $\ddot\gamma(t)$.
\end{TLCineApp}

\section{Sources}
\label{sec:cine-sources}

Ce cours reprend le document original, les systèmes pédagogiques
associés et des ressources de collègues de l'UPSTI. Les exercices sont
issus de sujets de concours ou d'activités explicitement mentionnés dans
le document source.


\clearpage
% ============================================================================
% COURS 04 — EXERCICES DE CINÉMATIQUE — PARTIE A
% Source : 5-Cinematique2022.docx
% ============================================================================

\section{Exercices du chapitre}
\label{sec:cine-exercices}

\subsection{Lecture d'une loi en trapèze}

Le profil d'accélération d'un mobile est défini par morceaux entre
$0$ et $3{,}2\ \mathrm s$. La position et la vitesse initiales sont
nulles.

\TLCineImage[.76\linewidth]{images/image102.png}

\begin{enumerate}
  \item Tracer la vitesse en intégrant graphiquement l'accélération.
  \item Tracer la position en calculant les aires sous la courbe de vitesse.
  \item Indiquer les valeurs maximales de vitesse et de déplacement.
  \item Vérifier la continuité de $x(t)$ et $v(t)$ aux changements de phase.
\end{enumerate}

\subsection{Maison DÔME}

Une maison circulaire de diamètre $16\ \mathrm m$ s'oriente afin de
suivre ou d'éviter le soleil. Lors de la coupure du moteur, sa vitesse
de rotation décroît linéairement jusqu'à zéro en $T_a=8\ \mathrm s$.
La fréquence initiale vaut $30\times10^{-3}\ \mathrm{tr\,min^{-1}}$.

\TLCineImage[.48\linewidth]{images/image103.pdf}

\begin{enumerate}
  \item Expliquer pourquoi la maison continue de tourner après la coupure.
  \item Calculer l'accélération angulaire en $\mathrm{rad\,s^{-2}}$.
  \item Déterminer l'angle parcouru pendant l'arrêt.
  \item Calculer la distance parcourue au niveau de la porte, située à $8\ \mathrm m$ de l'axe.
  \item Vérifier le critère de distance d'arrêt inférieure à $15\ \mathrm{cm}$.
\end{enumerate}

\subsection{Véhicule électrique F-City}

La F-City atteint $30\ \mathrm{km\,h^{-1}}$ en $5{,}5\ \mathrm s$, sa
vitesse maximale vaut $65\ \mathrm{km\,h^{-1}}$ et son autonomie
minimale est $80\ \mathrm{km}$.

\TLCineImage[.52\linewidth]{images/image104.pdf}

\begin{enumerate}
  \item Calculer l'accélération moyenne entre 0 et $30\ \mathrm{km\,h^{-1}}$.
  \item Calculer la durée maximale d'un trajet de $80\ \mathrm{km}$ à vitesse maximale.
  \item À partir du profil de vitesse fourni, déterminer les distances de chaque phase.
  \item Calculer la distance totale du parcours et vérifier l'autonomie.
\end{enumerate}

\subsection{Éolienne à pales orientables}

Le rotor 1 tourne autour de $(O,\vec x_0)$ avec l'angle $\theta$. Une
pale 2 pivote par rapport au rotor avec l'angle de calage $\alpha$.
On donne :
\[
\overrightarrow{OA}=a\vec y_1,
\qquad
\overrightarrow{AG}=x\vec x_2+y\vec y_2+z\vec z_2.
\]

\TLCineImage[.48\linewidth]{images/image105.jpeg}

\begin{enumerate}
  \item Tracer les deux figures de changement de base et donner les vecteurs rotation.
  \item Déterminer $\vec V(A\in1/0)$.
  \item Déterminer $\vec V(G\in2/0)$ par composition.
  \item Déterminer $\vec A(A\in1/0)$ puis $\vec A(G\in2/0)$ par dérivation.
  \item Identifier les termes centrifuge, tangentiel et de Coriolis.
\end{enumerate}

\subsection{Centrifugeuse d'entraînement}

Le bras 1 tourne autour de l'axe vertical $(O,\vec z)$ avec l'angle
$\alpha$. La cabine 2 pivote autour de $(A,\vec y_1)$ avec l'angle
$\beta$. On donne $\overrightarrow{OA}=a\vec x_1$ et
$\overrightarrow{AG}=b\vec z_2$.

\TLCineImage[.58\linewidth]{images/image106.png}

\begin{enumerate}
  \item Tracer les figures de changement de base et les vecteurs rotation.
  \item Déterminer $\vec\Omega_{2/0}$.
  \item Déterminer $\vec V(G\in2/0)$ par composition.
  \item Déterminer $\vec A(G\in2/0)$ par dérivation du résultat précédent.
  \item Interpréter les composantes de l'accélération ressentie par le pilote.
\end{enumerate}

\subsection{Robot manipulateur à deux axes}

Le corps 2 tourne par rapport au socle 1 autour de
$(O_1,\vec z_1)$ avec l'angle $\varphi$. Le bras 3 coulisse suivant
$\vec y_2=\vec y_3$ avec
$\overrightarrow{AG_3}=y(t)\vec y_3$ et
$\overrightarrow{O_1A}=a\vec x_2$.

\TLCineImage[.50\linewidth]{images/image109.jpeg}

\begin{enumerate}
  \item Donner la figure de changement de base et $\vec\Omega_{2/1}$.
  \item Écrire le torseur de la glissière 3/2.
  \item Déterminer $\vec V(G_3\in3/1)$.
  \item Déterminer $\vec A(G_3\in3/1)$.
  \item Repérer le terme de Coriolis lorsque $\dot y\neq0$ et $\dot\varphi\neq0$.
\end{enumerate}

\subsection{Transporteur horizontal}

Un chariot 2 se déplace suivant $\vec x_1$ avec la coordonnée $x(t)$.
Un bras 3 de longueur $L_3$ est articulé en $B$ sur le chariot ; son
centre de gravité est tel que
$\overrightarrow{BG_3}=L_3\vec y_3$ et son orientation est $\theta$.

\TLCineImage[.50\linewidth]{images/image110.jpeg}

\begin{enumerate}
  \item Donner les torseurs $\{\mathcal V_{2/1}\}_A$ et $\{\mathcal V_{3/2}\}_A$.
  \item Déterminer $\vec A(A\in2/1)$.
  \item Déterminer $\vec V(G_3\in3/1)$ puis $\vec A(G_3\in3/1)$.
  \item Exprimer les résultats en fonction de $\dot x$, $\ddot x$, $\dot\theta$, $\ddot\theta$ et $L_3$.
\end{enumerate}

\subsection{Centre d'usinage cinq axes HSM 600U}

Le centre possède trois translations $x(t)$, $y(t)$, $z(t)$ et deux
rotations $\alpha(t)$ et $\beta(t)$. Les axes linéaires atteignent
$40\ \mathrm{m\,min^{-1}}$ ; les axes rotatifs atteignent respectivement
$150$ et $250\ \mathrm{tr\,min^{-1}}$.

\begin{center}
\TLCineImage[.29\linewidth]{images/image111.jpeg}\hfill
\TLCineImage[.31\linewidth]{images/image112.png}\hfill
\TLCineImage[.34\linewidth]{images/image113.png}
\end{center}

\begin{enumerate}
  \item Identifier la nature des cinq mouvements d'axes.
  \item Déterminer les relations entre les bases des solides en translation.
  \item Écrire le vecteur position de l'extrémité d'outil $P$ dans $R_0$.
  \item Déterminer $\vec V(P\in5/0)$ et sa norme maximale.
  \item Déterminer la position, la vitesse et l'accélération du point $O_3$ porté par la table orientable.
  \item Retrouver la vitesse par composition des torseurs cinématiques.
  \item Écrire la vitesse relative de l'outil par rapport à un point de la pièce usinée.
\end{enumerate}

\subsection{Échelle pivotante automatique EPAS}

Le berceau 2 tourne par rapport à la tourelle 1 avec l'angle $\theta$.
L'échelle 3 se déploie suivant $\vec x_2$ avec la course $\lambda(t)$.
Pendant le dressage, la tourelle reste fixe.

\TLCineImage[.65\linewidth]{images/image172.png}

\begin{enumerate}
  \item Donner la nature des mouvements 2/0, 3/2 et 3/0.
  \item Décrire les trajectoires d'un point de l'échelle dans ces trois mouvements.
  \item Déterminer $\vec\Omega_{3/0}$.
  \item Déterminer la vitesse d'un point $G$ de la nacelle.
  \item Déterminer son accélération et identifier les termes liés à $\lambda$, $\theta$ et leurs dérivées.
  \item Vérifier le critère d'accélération maximale du cahier des charges.
\end{enumerate}

% ============================================================================
% COURS 04 — EXERCICES DE CINÉMATIQUE — PARTIE B
% Source : 5-Cinematique2022.docx
% ============================================================================

\subsection{Manipulateur de stèles de marbre}

\textit{Source : Jacques Le Goff.}

Le manipulateur permet l'élévation, l'inclinaison et l'orientation d'une
stèle pouvant atteindre $400\ \mathrm{kg}$. On étudie une combinaison
des mouvements d'élévation et d'inclinaison.

\begin{center}
\TLCineImage[.31\linewidth]{images/image197.pdf}\hfill
\TLCineImage[.31\linewidth]{images/image198.png}\hfill
\TLCineImage[.31\linewidth]{images/image200.png}
\end{center}

\begin{enumerate}
  \item Construire la chaîne de solides et les figures de changement de base.
  \item Déterminer, par dérivation du vecteur position, la vitesse du point $O_8$ par rapport à $R_0$.
  \item Retrouver le résultat par composition des vitesses.
  \item Déterminer l'accélération du point $O_8$.
  \item Identifier les termes dus à l'élévation, à l'inclinaison et à leur couplage.
\end{enumerate}

\subsection{Robot de manutention à parallélogramme}

Le robot déplace son poignet dans un plan à l'aide de deux moteurs. Les
solides 1, 2, 3 et 4 forment une structure en parallélogramme
déformable.

\begin{center}
\TLCineImage[.46\linewidth]{images/image222.png}\hfill
\TLCineImage[.46\linewidth]{images/image223.jpeg}
\end{center}

\begin{enumerate}
  \item Établir les relations entre les bases imposées par le parallélogramme.
  \item Déterminer les torseurs cinématiques des pivots commandés.
  \item Composer les torseurs pour obtenir les mouvements 3/0 et 4/0.
  \item Déterminer la vitesse d'un point du poignet par le champ des vitesses.
  \item Décrire sa trajectoire lorsque chacun des moteurs est successivement bloqué.
\end{enumerate}

\subsection{Monte-charge}

Une charge est entraînée par un câble et une poulie. Un tour moteur
correspond à une montée de $10\ \mathrm{cm}$. Le cycle comporte une
accélération jusqu'à un capteur situé à $0{,}5\ \mathrm m$, un palier à
$0{,}2\ \mathrm{m\,s^{-1}}$ jusqu'à $1{,}5\ \mathrm m$, puis un
freinage. La montée dure $15\ \mathrm s$.

\TLCineImage[.48\linewidth]{images/image238.png}

\begin{enumerate}
  \item Déterminer l'accélération de la première phase et les équations horaires.
  \item Calculer sa durée.
  \item Déterminer les équations horaires et la durée de la phase uniforme.
  \item Déterminer la décélération de la dernière phase.
  \item Calculer la distance totale de montée et tracer la vitesse moteur.
\end{enumerate}

\subsection{Robot de soudage quatre axes}

Le robot réalise un cordon linéaire puis un dégagement rapide de
l'organe terminal. Chaque axe possède son propre actionneur.

\begin{center}
\TLCineImage[.31\linewidth]{images/image239.jpeg}\hfill
\TLCineImage[.31\linewidth]{images/image241.png}\hfill
\TLCineImage[.31\linewidth]{images/image242.png}
\end{center}

\begin{enumerate}
  \item Déterminer les vitesses des points caractéristiques des solides 1 et 2.
  \item Donner le torseur du solide 2 par rapport au bâti, réduit en $O_2$.
  \item Déterminer la vitesse de l'extrémité de soudage.
  \item Établir les relations entre les paramètres pour obtenir une trajectoire rectiligne imposée.
  \item Lorsque deux axes sont bloqués, déterminer l'accélération pendant le dégagement rapide.
\end{enumerate}

\subsection{Scénic contre Ferrari de Formule 1}

Les deux véhicules effectuent : accélération de 0 à
$100\ \mathrm{km\,h^{-1}}$, maintien pendant $10\ \mathrm s$, puis
freinage jusqu'à l'arrêt. Le Scénic atteint $100\ \mathrm{km\,h^{-1}}$
en $13{,}4\ \mathrm s$ et freine sur $46\ \mathrm m$ ; la F1 atteint
la même vitesse en $2{,}3\ \mathrm s$ et freine sur $14\ \mathrm m$.

\begin{center}
\TLCineImage[.45\linewidth]{images/image254.jpeg}\hfill
\TLCineImage[.45\linewidth]{images/image255.jpeg}
\end{center}

\begin{enumerate}
  \item Établir $a(t)$, $v(t)$ et $x(t)$ pour les trois phases de chaque véhicule.
  \item Calculer la distance totale parcourue.
  \item Calculer le temps et la distance nécessaires pour atteindre la vitesse maximale de chaque véhicule.
  \item Comparer les accélérations et décélérations moyennes.
\end{enumerate}

\subsection{Manipulateur de collecteur d'échappement}

\textit{Source : Mines PTSI 2006.}

Le manipulateur assiste un opérateur pour déplacer un collecteur de
$17\ \mathrm{kg}$. Un vérin équilibre le poids ; les autres mouvements
sont guidés manuellement.

\begin{center}
\TLCineImage[.34\linewidth]{images/image256.jpeg}\hfill
\TLCineImage[.28\linewidth]{images/image257.png}\hfill
\TLCineImage[.30\linewidth]{images/image258.png}
\end{center}

\begin{enumerate}
  \item Compléter le schéma cinématique et le paramétrage.
  \item Déterminer les vecteurs rotation associés aux différentes liaisons.
  \item Écrire les torseurs cinématiques au point de réduction adapté.
  \item Déterminer la vitesse d'un point $P$ du collecteur par composition.
  \item Discuter les positions singulières et les directions de vitesse accessibles.
\end{enumerate}

\subsection{Robot de pick and place}

Le schéma cinématique d'un robot de prise-dépose comporte plusieurs
pivots et une glissière.

\begin{center}
\TLCineImage[.43\linewidth]{images/image264.png}\hfill
\TLCineImage[.43\linewidth]{images/image265.png}
\end{center}

\begin{enumerate}
  \item Déterminer les vecteurs rotation des solides successifs.
  \item Écrire les torseurs cinématiques des liaisons.
  \item Déterminer les vitesses de points par déplacement des torseurs.
  \item Construire les figures de changement de base.
  \item Déterminer la vitesse et l'accélération de l'effecteur.
  \item Identifier les mouvements de translation possibles entre deux solides.
\end{enumerate}

\subsection{Shimmy d'une roue de train d'atterrissage}

L'avion 1 est en translation rectiligne uniforme par rapport au sol 0.
L'étrier 2 pivote autour d'un axe vertical et porte la roue 3, de rayon
$R$, roulant sur la piste au point $A$. Le mouvement de shimmy est une
oscillation de l'étrier autour de son pivot.

\TLCineImage[.66\linewidth]{images/image287.png}

\begin{enumerate}
  \item Écrire la condition de roulement sans glissement en $A$.
  \item Appliquer la composition des vitesses pour déterminer la vitesse du centre $O_3$.
  \item Relier la vitesse du centre et celle du point de contact par le champ des vitesses.
  \item Déduire la relation entre la vitesse de l'avion, la rotation de la roue et l'oscillation de l'étrier.
\end{enumerate}

\subsection{Canne robotisée}

Une roue motorisée de rayon $R$ déplace une canne inclinée d'un angle
$\theta(t)$. La longueur $AH=l(t)$ varie afin de maintenir la poignée
$H$ à la hauteur constante $h_0$.

\begin{center}
\TLCineImage[.43\linewidth]{images/image293.png}\hfill
\TLCineImage[.43\linewidth]{images/image294.png}
\end{center}

\begin{enumerate}
  \item Écrire la fermeture géométrique $\overrightarrow{IA}+\overrightarrow{AH}+\overrightarrow{HI}=\vec0$.
  \item Projeter pour relier $l(t)$, $\theta(t)$, $R$ et $h_0$.
  \item Écrire la condition de roulement sans glissement en $I$.
  \item Déterminer $\vec V(H\in3/2)$ et les vecteurs rotation.
  \item Composer les vitesses pour obtenir $\vec V(H\in3/0)$.
  \item Déduire la vitesse de la roue assurant la vitesse de déplacement imposée sans mouvement parasite de la main.
\end{enumerate}

\subsection{Dépose de bagage automatisée}

Le système déplace et bascule un bagage entre différents convoyeurs. La
cinématique doit garantir la continuité des vitesses et limiter les
accélérations transmises au bagage.

\begin{center}
\TLCineImage[.29\linewidth]{images/image295.png}\hfill
\TLCineImage[.29\linewidth]{images/image296.png}\hfill
\TLCineImage[.34\linewidth]{images/image297.png}
\end{center}

\begin{enumerate}
  \item Identifier les mouvements élémentaires et les trajectoires des points du bagage.
  \item Écrire les torseurs des liaisons puis les composer.
  \item Déterminer la vitesse et l'accélération d'un point caractéristique.
  \item Vérifier les contraintes de temps de transfert et de maintien du bagage.
\end{enumerate}

\subsection{Véhicule autoguidé hospitalier}

Le VAG doit s'arrêter depuis $v_{\max}=1{,}2\ \mathrm{m\,s^{-1}}$ sur
une distance maximale $d_f=0{,}2\ \mathrm m$. La vitesse décroît
linéairement pendant le freinage.

\begin{center}
\TLCineImage[.40\linewidth]{images/image298.pdf}\hfill
\TLCineImage[.48\linewidth]{images/image299.png}
\end{center}

\begin{enumerate}
  \item Exprimer la durée du freinage en fonction de $d_f$ et $v_{\max}$.
  \item Calculer la décélération constante nécessaire.
  \item Tracer les courbes $v(t)$, $a(t)$ et $x(t)$.
  \item Vérifier la compatibilité avec la condition de non-dérapage.
\end{enumerate}

\subsection{Synthèse : choix d'une méthode}

Pour chacun des cas suivants, indiquer la méthode la plus efficace avant
de réaliser le calcul demandé :
\begin{enumerate}
  \item vitesse de deux points d'un même solide ;
  \item vitesse d'un point dans une chaîne de trois solides ;
  \item dérivée d'un vecteur unitaire d'une base mobile ;
  \item relation cinématique imposée par un contact sans glissement ;
  \item déplacement associé à une loi de vitesse en trapèze ;
  \item accélération d'un point déjà connu par son vecteur vitesse.
\end{enumerate}


\printindex
